\documentclass[11pt]{article}
\usepackage[margin=1.05in]{geometry}
\usepackage{amsmath,amssymb,amsthm,booktabs,microtype,xcolor}
\usepackage[colorlinks=true,linkcolor=blue!50!black,urlcolor=blue!50!black]{hyperref}
\newtheorem{theorem}{Theorem}
\newtheorem{lemma}[theorem]{Lemma}
\newtheorem{corollary}[theorem]{Corollary}
\newcommand{\R}{\mathbb R}
\newcommand{\N}{\mathbb N}
\setlength{\emergencystretch}{2em}
\title{Bounded Collatz Correction and Reciprocal Orbit Sums}
\author{Collatz proof project\thanks{Prepared with Codex (OpenAI). All central statements are formalized in Lean. No claim of literature priority is made.}}
\date{September 5, 2026}
\begin{document}
\maketitle
\begin{abstract}
For a positive integer shortcut Collatz orbit, boundedness of the normalized affine correction is equivalent to convergence of the sum of reciprocal orbit values. Two finite inequalities prove the equivalence: a telescoping estimate that charges reciprocal mass to correction increments, and an exponential upper bound on correction in terms of accumulated reciprocal mass. No strict odd-step density gap is assumed. Bounded correction consequently forces escape from every finite set. The equivalence characterizes a class of hypothetical counterexamples; it neither excludes that class nor shows that every escaping orbit belongs to it.
\end{abstract}

\section{Definitions and exact characterization}
Let $N\in\N$ be positive and define
\[
T(x)=\begin{cases}x/2&x\text{ even},\\(3x+1)/2&x\text{ odd},\end{cases}
\qquad n_t=T^t(N).
\]
All $n_t$ are positive integers. Write $j_t$ for the number of odd terms among $n_0,\ldots,n_{t-1}$, and let $d_t$ be the integer affine correction determined by
\[
2^t n_t=3^{j_t}N+d_t.
\]
Set
\[
c_t=\frac{d_t}{3^{j_t}},\qquad r_t=\frac1{n_t},\qquad
S_L=\sum_{t=0}^{L-1}r_t.
\]
Thus $c_0=0$, $c_t\ge0$, and $c_t$ is nondecreasing. The existing Lean predicate \texttt{Supercritical N} means precisely $\sup_t c_t<\infty$. Its name does not add any density assumption.

\begin{theorem}\label{thm:main}
For every positive integer $N$,
\[
\sup_t c_t<\infty\quad\Longleftrightarrow\quad
\sum_{t=0}^{\infty}\frac1{n_t}<\infty.
\]
\end{theorem}
The proof below uses finite estimates valid for every positive seed. In particular, neither direction assumes that the orbit diverges, has a limiting odd-step density, or satisfies a uniform density gap.

\newpage
\section{Finite inequalities}
The affine identity gives
\begin{equation}\label{eq:inverse}
\frac{2^t}{3^{j_t}}=(N+c_t)r_t.
\end{equation}
The correction cocycle gives
\begin{equation}\label{eq:increment}
c_{t+1}-c_t=
\begin{cases}
0&n_t\text{ even},\\
(N+c_t)r_t/3&n_t\text{ odd}.
\end{cases}
\end{equation}
These identities hold in $\R$; the divisions are legitimate because $n_t>0$.

\begin{lemma}\label{lem:ledger}
For every $L\in\N$,
\begin{equation}\label{eq:ledger}
N S_L\le6c_L+N(r_L-r_0).
\end{equation}
\end{lemma}
\begin{proof}
We prove the one-step estimate
\[
Nr_t\le6(c_{t+1}-c_t)+N(r_{t+1}-r_t).
\]
If $n_t$ is even, then $r_{t+1}=2r_t$ and the correction increment vanishes, so equality holds. If $n_t$ is odd, equation~\eqref{eq:increment} and $c_t\ge0$ give
\[
6(c_{t+1}-c_t)\ge2Nr_t.
\]
Adding $N(r_{t+1}-r_t)$ gives at least $Nr_t$, because $r_{t+1}\ge0$. Sum over $0\le t<L$ and telescope, using $c_0=0$.
\end{proof}
The constant six is a convenient uniform choice; optimality is not asserted.

\begin{lemma}\label{lem:exp}
For every $L\in\N$,
\begin{equation}\label{eq:exp}
N+c_L\le N\exp(S_L).
\end{equation}
\end{lemma}
\begin{proof}
In both parity cases equation~\eqref{eq:increment} implies
\[
N+c_{t+1}\le(N+c_t)(1+r_t)\le(N+c_t)e^{r_t}.
\]
All factors are nonnegative. Induction from $c_0=0$, using $e^{a+b}=e^a e^b$, proves the claim. The Lean proof uses the existing one-step bound $\texttt{idealC}\ 1\ n_t\le1$ and the cocycle; the sharper odd-step constant $1/3$ is unnecessary here.
\end{proof}

\newpage
\section{Infinite sums and escape}
\begin{proof}[Proof of Theorem~\ref{thm:main}]
Suppose first that $c_t\le B$ for every $t$. Positivity and integrality give $0\le r_t\le1$. Lemma~\ref{lem:ledger} therefore implies
\[
S_L\le\frac{6B}{N}+1\qquad\text{for every }L.
\]
The nonnegative partial sums are bounded, so the series converges. Lean uses mathlib's \texttt{summable\_of\_sum\_range\_le}.

Conversely, suppose $S=\sum_{t\ge0}r_t<\infty$. Nonnegativity gives $S_L\le S$. Lemma~\ref{lem:exp} yields
\[
N+c_L\le Ne^S\qquad\text{for every }L,
\]
which bounds $c_L$. The formal proof uses the upper bound $Ne^S$ for $c_L$; the displayed estimate also gives $c_L\le N(e^S-1)$.
\end{proof}

\begin{corollary}\label{cor:escape}
If $N>0$ and $\sup_t c_t<\infty$, then $n_t\to+\infty$. More precisely, for every $B\in\N$ there is a time after which $n_t\ge B$ always holds.
\end{corollary}
\begin{proof}
The terms of a convergent series tend to zero, so $r_t\to0$. Eventually $r_t<1/(B+1)$, which forces $n_t>B$. The Lean conclusion uses the equivalent non-strict at-top formulation.
\end{proof}
In particular, a positive orbit that is eventually periodic cannot have bounded correction. This last sentence follows immediately from Corollary~\ref{cor:escape}; it is not a separately named declaration in the new module.

\section{What this changes in the proof program}
The earlier geometric inverse-drift bound remains a sufficient condition for bounded correction. Theorem~\ref{thm:main} supplies an exact orbit-value condition instead, and does not require exponential growth. For example, an independently proved envelope $n_t\ge A(t+1)^2$ with $A>0$ would suffice by comparison with a convergent square reciprocal series. This comparison is an explanatory consequence, not a separate formal theorem in this module, and no such envelope is claimed for any natural Collatz seed.

The main unresolved question is now visible in orbit terms: can an escaping positive integer Collatz orbit have $\sum_t1/n_t=\infty$? Mere escape does not imply reciprocal summability for general positive sequences. Even settling this coverage question would leave the task of excluding bounded-correction orbits, as well as the separate nontrivial-cycle problem. Thus the present result is a characterization, not a proof of Collatz or an exclusion theorem.

\newpage
\section{Lean declarations and reproducibility}
The implementation is \texttt{lean/Collatz/Reciprocal.lean}. It imports the project's ideal-correction bounds and mathlib's real exponential and nonnegative-series facts. The primary theorem is:
\begin{quote}\small\ttfamily
supercritical\_iff\_summable\_reciprocal\par
\normalfont For $N>0$, \texttt{Supercritical N} if and only if\par
\texttt{Summable (orbitReciprocal N)}.
\end{quote}
Here \texttt{orbitReciprocal N t} is $1/(T^t(N):\R)$.

\begin{center}
\small
\begin{tabular}{lp{9.4cm}}
\toprule
Result&Declaration in namespace \texttt{Collatz}\\
\midrule
Equation~\eqref{eq:inverse}&\texttt{inverse\_drift\_eq\_reciprocal}\\
Lemma~\ref{lem:ledger}&\texttt{reciprocal\_sum\_ledger}\\
Lemma~\ref{lem:exp}&\texttt{idealC\_exp\_reciprocal\_bound}\\
Theorem~\ref{thm:main}&\texttt{supercritical\_iff\_summable\_reciprocal}\\
Corollary~\ref{cor:escape}&\texttt{supercritical\_orbit\_tendsto\_atTop}\\
\bottomrule
\end{tabular}
\end{center}
The module also contains each implication separately and the local one-step ledger. No \texttt{sorry}, new axiom, or \texttt{native\_decide} is used.

From the repository root, rebuild and audit with:
\begin{verbatim}
cd lean
lake build Collatz.Reciprocal
cd ..
python3 analysis/audit_theorems.py --build
\end{verbatim}
The pinned Lean toolchain is \texttt{leanprover/lean4:v4.27.0-rc1}. The selected audit now contains 48 declarations, including the two finite estimates, the equivalence, and the escape corollary. Their reported dependencies use only \texttt{propext}, \texttt{Classical.choice}, and \texttt{Quot.sound}. The JSON audit records exact printed types and repository provenance. It checks selected modules and their Lean-reported dependencies; it is not an independent kernel replay or a rebuild of every nested library.

\section{Sources and scope of attribution}
The definitions and affine cocycle are in the existing project module \texttt{Collatz.Ideal}; the previously established geometric criterion is in \texttt{Collatz.IdealBounds}. Nonnegative-series and exponential lemmas are supplied by the pinned mathlib checkout. All new mathematical arguments used here are displayed above and implemented in \texttt{Collatz.Reciprocal}. No external literature-priority claim is made.
\end{document}
